\documentclass[a4paper,12pt]{article}
\usepackage[T2A]{fontenc}
\usepackage[utf8]{inputenc}
\usepackage[russian]{babel}
\usepackage{amsmath,amssymb,mathtools}
\usepackage{helvet}
\renewcommand{\familydefault}{\sfdefault}
\usepackage{geometry}
\geometry{left=23mm,right=22mm,top=23mm,bottom=22mm}
\usepackage{microtype}
\usepackage{icomma}
\usepackage{xcolor}
\usepackage{tikz}
\usetikzlibrary{arrows.meta,positioning,shapes.geometric,calc}
\usepackage{tkz-euclide}
\usepackage{tabularx,booktabs,array}
\usepackage{enumitem}
\usepackage{fancyhdr}
\usepackage{setspace}
\usepackage{hyperref}
\definecolor{accent}{HTML}{087C82}
\definecolor{darkaccent}{HTML}{07545B}
\definecolor{textgray}{HTML}{323B3D}
\definecolor{muted}{HTML}{647074}
\hypersetup{colorlinks=true,linkcolor=darkaccent,urlcolor=darkaccent,pdfauthor={Лёвин Артём Александрович},pdftitle={Чек-лист. Подобие и площади в трапеции. 20.11.2026}}
\setlength{\parindent}{0pt}
\setlength{\parskip}{4pt}
\setstretch{1.21}
\setlist[itemize]{leftmargin=*,itemsep=2pt,topsep=3pt}
\setlist[enumerate]{leftmargin=*,itemsep=8pt,topsep=5pt}
\pagestyle{fancy}
\fancyhf{}
\fancyhead[L]{\footnotesize\color{muted}Площади и подобие в трапеции}
\fancyhead[R]{\footnotesize\color{muted}София Кальней}
\fancyfoot[C]{\footnotesize\thepage}
\renewcommand{\headrulewidth}{0.25pt}
\setlength{\headheight}{15pt}
\newcommand{\key}[1]{\textcolor{darkaccent}{\textbf{#1}}}
\newcommand{\met}[1]{\textcolor{accent}{#1}}
\newcommand{\Area}{S}
\newcommand{\smallrule}{\par\vspace{1mm}\noindent{\color{accent}\rule{\linewidth}{.45pt}}\par\vspace{1mm}}
\begin{document}
\thispagestyle{empty}
\begin{center}
\vspace*{26mm}
{\Large\color{darkaccent} МАТЕМАТИКА}\par
\vspace{16mm}
{\Huge\color{textgray} Чек-лист}\par
\vspace{12mm}
{\LARGE\color{accent} Подобие треугольников\\[3mm] и площади в трапеции}\par
\vspace{12mm}
{\large Геометрия: задачи повышенной сложности}\par
\vspace{32mm}
{\large София Кальней}\par
\vspace{10mm}
{\large 20 ноября 2026 г.}\par
\vfill
{\small Лёвин Артём Александрович\\эксклюзивно для Софии Кальней}\par
\vspace{14mm}
\begin{tikzpicture}[x=1cm,y=1cm]
\draw[accent,thin] (0,0) .. controls (1.3,.52) and (2.3,-.55) .. (4,0)
                     .. controls (5.4,.53) and (6.1,-.55) .. (8,0);
\fill[accent] (4,0) circle (0.9pt);
\end{tikzpicture}
\end{center}
\clearpage

\section*{1. Главная идея урока}
В задачах на площади сложных фигур часто достаточно найти \key{отношения отрезков}. Для этого удобно выделить треугольники, обнаружить подобие и при необходимости сделать дополнительное построение.

\smallrule
\key{Маршрут:} нужная площадь $\to$ опорный треугольник $\to$\\
отношение площадей $\to$ подобие / средняя линия $\to$ ответ.

\section*{2. Формулы, которые нужно узнавать сразу}
\begin{tabularx}{\linewidth}{@{}>{\raggedright\arraybackslash}p{.38\linewidth} X@{}}
\toprule
\textbf{Ситуация} & \textbf{Формула и смысл}\\
\midrule
Треугольник & $\displaystyle S_\triangle=\frac12 ah$, где $h$ --- высота к стороне $a$.\\[1mm]
Трапеция & $\displaystyle S_{\text{тр}}=\frac{a+b}{2}H$, где $a,b$ --- основания, $H$ --- высота.\\[1mm]
Средняя линия трапеции & $\displaystyle m=\frac{a+b}{2}$; соединяет середины боковых сторон.\\[1mm]
Общая высота & $\displaystyle \frac{S_1}{S_2}=\frac{a_1}{a_2}$, если высоты к сравниваемым основаниям равны.\\[1mm]
Общий угол & $\displaystyle \frac{S_1}{S_2}=\frac{u_1v_1}{u_2v_2}$, если угол между соответствующими парами сторон один и тот же.\\[1mm]
Подобные треугольники & Соответственные стороны относятся как $k$, а площади --- как $k^2$.\\
\bottomrule
\end{tabularx}

\section*{3. Средняя линия и половина высоты}
Если $ABCD$ --- трапеция с основаниями $AD\parallel BC$, $P$ --- середина $AB$, $M$ --- середина $CD$, то
\[
PM\parallel AD\parallel BC,\qquad PM=\frac{AD+BC}{2}.
\]
Точка $P$ находится на половине высоты трапеции. Значит, высота треугольника $APD$, проведённая к $AD$, равна $H/2$:
\[
S_{APD}=\frac12\,AD\cdot\frac H2=\frac{AD\cdot H}{4},\qquad
\boxed{\frac{S_{APD}}{S_{ABCD}}=\frac{AD}{2(AD+BC)}}.
\]
То же утверждение верно для $M$ и треугольника $AMD$.

\section*{4. Как найти подобные треугольники}
\begin{itemize}
\item \key{Вертикальные углы} дают одну пару равных углов.
\item \key{Параллельные прямые} дают вторую пару равных углов (накрест лежащие или соответственные).
\item \key{Два равных угла} доказывают подобие треугольников; затем записывается пропорция \emph{соответственных} сторон.
\end{itemize}
\textcolor{muted}{Важно: только равенства вертикальных углов для доказательства подобия недостаточно.}

\subsection*{Модель: две параллельные стороны и пересекающиеся прямые}
Пусть $UV\parallel WZ$, а прямые $UZ$ и $VW$ пересекаются в точке $O$.

\begin{center}
\begin{tikzpicture}[scale=.88]
  \tkzDefPoint(0,0){U}
  \tkzDefPoint(0,4){V}
  \tkzDefPoint(5.3333333,2){O}
  \tkzDefPoint(8,1){W}
  \tkzDefPoint(8,3){Z}
  \tkzDrawSegments[thin,color=darkaccent](U,V W,Z U,Z V,W)
  \tkzDrawPoints[size=2.5](U,V,O,W,Z)
  \tkzLabelPoint[below](U){$U$}
  \tkzLabelPoint[above](V){$V$}
  \tkzLabelPoint[below](O){$O$}
  \tkzLabelPoint[below](W){$W$}
  \tkzLabelPoint[above](Z){$Z$}
\end{tikzpicture}
\end{center}

\begin{itemize}
\item $\angle UOV=\angle ZOW$ --- вертикальные углы.
\item $\angle UVO=\angle ZWO$ --- углы при параллельных $UV$ и $WZ$.
\end{itemize}
Следовательно, $\triangle UVO\sim\triangle ZWO$ по двум углам. Порядок вершин показывает соответствия $U\leftrightarrow Z$, $V\leftrightarrow W$, $O\leftrightarrow O$:
\[
\boxed{\frac{UV}{ZW}=\frac{UO}{ZO}=\frac{VO}{WO}}.
\]
\key{Практический вывод.} Увидев пересечение прямых, найдите вторую пару равных углов: часто она появляется после проведения параллельной прямой или продолжения стороны трапеции.
\clearpage

\section*{5. Пример 1. Продление стороны и подобие}
\textbf{Задано.} Площадь трапеции $ABCD$ равна $30$, $AD\parallel BC$, $AD=2BC$. Точка $P$ --- середина $AB$, а точка $R$ на стороне $CD$ выбрана так, что $2CD=3RD$. Прямые $AR$ и $PD$ пересекаются в точке $Q$. Найти $S_{APQ}$.

\begin{center}
\begin{tikzpicture}[scale=.87,>=Stealth]
  \tkzDefPoint(0,0){A}
  \tkzDefPoint(7,0){D}
  \tkzDefPoint(.5,4){B}
  \tkzDefPoint(4,4){C}
  \tkzDefPoint(.25,2){P}
  \tkzDefPoint(5,2.6666667){R}
  \tkzDefPoint(2.5,1.3333333){Q}
  \tkzDefPoint(7.5,4){T}
  \tkzDrawPolygon[thin](A,B,C,D)
  \tkzDrawSegments[thin,color=accent](C,T A,T P,D D,T)
  \tkzDrawPoints[size=2.4](A,B,C,D,P,R,Q,T)
  \tkzLabelPoint[below left](A){$A$}
  \tkzLabelPoint[above left](B){$B$}
  \tkzLabelPoint[above](C){$C$}
  \tkzLabelPoint[below](D){$D$}
  \tkzLabelPoint[left](P){$P$}
  \tkzLabelPoint[right](R){$R$}
  \tkzLabelPoint[below right](Q){$Q$}
  \tkzLabelPoint[above](T){$T$}
\end{tikzpicture}
\end{center}

\textbf{Шаг 1.} Так как $CD=CR+RD$ и $2CD=3RD$, то
\[
2(CR+RD)=3RD\quad\Longrightarrow\quad CR:RD=1:2.
\]
\textbf{Шаг 2.} Продлим $BC$ до пересечения с $AR$ в точке $T$. Из подобия $\triangle RCT\sim\triangle RDA$ (вертикальные углы при $R$, параллельные основания) следует
\[
\frac{CT}{AD}=\frac{CR}{RD}=\frac12.
\]
Поэтому $CT=AD/2=BC$, откуда $BT=BC+CT=AD$. Четырёхугольник $ABTD$ --- параллелограмм, поэтому $AP\parallel DT$ и $DT=AB=2AP$.

\textbf{Шаг 3.} Треугольники $APQ$ и $TDQ$ подобны по двум углам:
\[
\frac{PQ}{QD}=\frac{AP}{DT}=\frac12
\quad\Longrightarrow\quad
\boxed{\frac{PQ}{PD}=\frac13}.
\]
\textbf{Шаг 4.} Пусть $BC=b$, высота трапеции $H$. Тогда
\[
S_{APD}=\frac{2b\cdot H}{4}=\frac{bH}{2},\quad
S_{ABCD}=\frac{b+2b}{2}H=\frac{3bH}{2},\quad
S_{APD}=\frac13\cdot30=10.
\]
Треугольники $APQ$ и $APD$ имеют общую высоту к прямой $PD$, так что
\[
\boxed{S_{APQ}=\frac{PQ}{PD}\,S_{APD}=\frac13\cdot10=\frac{10}{3}}.
\]
\clearpage

\section*{6. Пример 2. Две средние линии и подобие}
\textbf{Задано.} Площадь трапеции $ABCD$ равна $64$, $AD\parallel BC$, $AD=3BC$. Точка $M$ --- середина $CD$, точка $N\in AB$ такова, что $AN:NB=1:3$. Прямые $AM$ и $DN$ пересекаются в точке $K$. Найти $S_{DMK}$.

\begin{center}
\begin{tikzpicture}[scale=.72]
  \tkzDefPoint(0,0){A}
  \tkzDefPoint(9,0){D}
  \tkzDefPoint(.8,4){B}
  \tkzDefPoint(3.8,4){C}
  \tkzDefPoint(.4,2){E}
  \tkzDefPoint(6.4,2){M}
  \tkzDefPoint(.2,1){N}
  \tkzDefPoint(3.2,1){X}
  \tkzDefPoint(2.4,.75){K}
  \tkzDrawPolygon[thin](A,B,C,D)
  \tkzDrawSegments[thin,color=accent](E,M A,M D,N N,X)
  \tkzDrawPoints[size=2.3](A,B,C,D,E,M,N,X,K)
  \tkzLabelPoint[below left](A){$A$}
  \tkzLabelPoint[above left](B){$B$}
  \tkzLabelPoint[above](C){$C$}
  \tkzLabelPoint[below](D){$D$}
  \tkzLabelPoint[left](E){$E$}
  \tkzLabelPoint[right](M){$M$}
  \tkzLabelPoint[left](N){$N$}
  \tkzLabelPoint[above](X){$X$}
  \tkzLabelPoint[below](K){$K$}
\end{tikzpicture}
\end{center}

\textbf{Шаг 1.} Отметим середину $E$ стороны $AB$. Отрезок $EM$ --- средняя линия трапеции. Если $BC=b$, то
\[
AD=3b,\qquad EM=\frac{b+3b}{2}=2b.
\]
\textbf{Шаг 2.} Из $AN:NB=1:3$ следует $AN=AB/4$. Поскольку $AE=AB/2$, точка $N$ --- середина $AE$. Проведём $NX\parallel EM$ до пересечения с $AM$. В треугольнике $AEM$ отрезок $NX$ --- средняя линия:
\[
AX=XM,\qquad NX=\frac{EM}{2}=b.
\]
\textbf{Шаг 3.} В треугольниках $NKX$ и $DKA$ есть вертикальные углы при $K$, а $NX\parallel AD$. Поэтому они подобны:
\[
\frac{KX}{KA}=\frac{NX}{AD}=\frac{b}{3b}=\frac13.
\]
Положим $KX=u$, тогда $AK=3u$, $AX=4u$, $XM=4u$,
\[
KM=KX+XM=5u,\qquad AM=8u,\qquad\boxed{\frac{KM}{AM}=\frac58}.
\]
\textbf{Шаг 4.} Высота треугольника $AMD$ вдвое меньше высоты трапеции:
\[
\frac{S_{AMD}}{S_{ABCD}}=\frac{3b}{2(3b+b)}=\frac38,
\qquad S_{AMD}=\frac38\cdot64=24.
\]
У треугольников $KMD$ и $AMD$ общая высота к прямой $AM$, поэтому
\[
\boxed{S_{KMD}=\frac{KM}{AM}\,S_{AMD}=\frac58\cdot24=15}.
\]
\clearpage

\section*{7. Памятка: как избежать типичных ошибок}
\begin{itemize}
\item \key{Проверяйте порядок букв.} Из $CR:RD=1:2$ следует $CD=3$ части, $RD=2$ части. Именно отрезок $CD$ состоит из двух частей $CR$ и $RD$.
\item \key{Различайте «часть от части» и «часть от целого».} Если $PQ:QD=1:2$, то $PQ:PD=1:3$.
\item \key{Сравнивайте стороны, лежащие друг против друга под равными углами.} Пропорция из подобия требует правильного соответствия вершин.
\item \key{Не переносите коэффициент подобия напрямую на площади.} Для самих подобных треугольников отношение площадей равно квадрату коэффициента; при общей высоте отношение площадей равно отношению оснований.
\item \key{Обосновывайте дополнительные построения.} Линия должна дать параллельность, среднюю линию или пару подобных треугольников.
\item \key{Проверяйте результат.} Внутренний треугольник, лежащий внутри трапеции, не может иметь площадь больше площади всей трапеции.
\end{itemize}
\smallrule
\textbf{Приём для самостоятельной работы.} Перед вычислениями запишите целевую дробь $S_{\text{мал}}/S_{\text{опор}}$ и решите, какие длины или высоты осталось найти. На чертеже удобно помечать соотношения как $u$, $2u$, $3u$ --- это уменьшает риск ошибиться в долях.

\clearpage
\thispagestyle{fancy}
\begin{center}{\LARGE\color{accent}\textbf{Алгоритм: площади и подобие в трапеции}}\end{center}
\vspace{8mm}
\begin{center}
\tikzset{
  fstart/.style={ellipse,draw=accent,line width=.75pt,minimum width=6.3cm,minimum height=.9cm,align=center,font=\normalsize},
  fstep/.style={rectangle,rounded corners=2pt,draw=muted,line width=.65pt,text width=9.5cm,minimum height=1.03cm,align=center,font=\normalsize,inner sep=6pt},
  fchoice/.style={diamond,draw=accent,line width=.8pt,aspect=3.2,text width=6.7cm,align=center,inner xsep=3pt,inner ysep=2pt,font=\normalsize},
  farrow/.style={-{Stealth[length=2.4mm]},line width=.75pt,color=darkaccent}
}
\begin{tikzpicture}[node distance=7mm and 15mm]
\node[fstart] (start) {Прочитайте условие и отметьте все отношения};
\node[fstep,below=of start] (s1) {Выберите треугольник, площадь которого нужно найти};
\node[fstep,below=of s1] (s2) {Сравните его площадь с площадью удобной опорной фигуры};
\node[fstep,below=of s2] (s3) {Найдите недостающие отношения: при необходимости продлите сторону или проведите среднюю линию};
\node[fstep,below=of s3] (s4) {Докажите подобие по двум углам; запишите пропорции соответствующих сторон};
\node[fchoice,below=11mm of s4] (dec) {Все длины и доли площади определены?};
\node[fstep,below=12mm of dec] (s5) {Вычислите нужную площадь через отношение к опорной фигуре};
\node[fstart,below=of s5] (end) {Проверьте доли, единицы и правдоподобие ответа};
\draw[farrow] (start)--(s1);
\draw[farrow] (s1)--(s2);
\draw[farrow] (s2)--(s3);
\draw[farrow] (s3)--(s4);
\draw[farrow] (s4)--(dec);
\draw[farrow] (dec)--node[right,font=\small]{да}(s5);
\draw[farrow] (s5)--(end);
\draw[farrow] (dec.west)--++(-1.9,0) |- node[pos=.1,above,font=\small]{нет} (s3.west);
\end{tikzpicture}
\end{center}
\clearpage

\section*{Тренировочный блок}
\textcolor{muted}{Путь А --- 10 задач. От базовых отношений к задачам с дополнительными построениями. Чертежи выполните самостоятельно; во всех задачах основания трапеции $ABCD$ --- $AD$ и $BC$, причём $AD\parallel BC$.}
\begin{enumerate}
\item \textbf{Базовый.} Основания трапеции равны $12$ и $8$, высота равна $5$. Найдите площадь трапеции.
\item \textbf{Базовый.} Основания трапеции равны $18$ и $10$. Точки $E$ и $M$ --- середины её боковых сторон. Найдите длину $EM$.
\item \textbf{Базовый.} $AD=16$, $BC=8$, высота трапеции $6$. Точка $P$ --- середина $AB$. Найдите площадь треугольника $APD$.
\item \textbf{Средний.} Площадь треугольника $ACD$ равна $50$. Точка $K$ лежит на $AC$, причём $AK:KC=2:3$. Найдите площадь треугольника $AKD$.
\item \textbf{Средний.} Площадь треугольника $APD$ равна $42$. Точка $Q$ лежит на отрезке $PD$, причём $PQ:QD=2:5$. Найдите площадь треугольника $APQ$.
\item \textbf{Средний.} $AD=3BC$, площадь трапеции равна $80$. Точка $M$ --- середина $CD$. Найдите площадь треугольника $AMD$.
\item \textbf{Выше среднего.} $AD=3BC$. Точка $P$ --- середина $AB$, точка $R$ на $CD$ удовлетворяет $3CD=5RD$. Прямые $AR$ и $PD$ пересекаются в точке $Q$. Найдите отношение $PQ:PD$.
\item \textbf{Выше среднего.} Площадь трапеции равна $90$, $AD=2BC$. Точка $P$ --- середина $AB$, $R\in CD$ и $2CD=3RD$. Прямые $AR$ и $PD$ пересекаются в точке $Q$. Найдите $S_{APQ}$.
\item \textbf{Сложный.} Площадь трапеции равна $128$, $AD=3BC$. Точка $M$ --- середина $CD$, а точка $N\in AB$ такова, что $AN:NB=1:3$. Прямые $AM$ и $DN$ пересекаются в точке $K$. Найдите $S_{DMK}$.
\item \textbf{Сложный.} Площадь трапеции равна $108$, $AD=2BC$. Точка $M$ --- середина $CD$, точка $N\in AB$ удовлетворяет $AN:NB=1:2$. Прямые $AM$ и $DN$ пересекаются в точке $K$. Найдите $S_{DMK}$.
\end{enumerate}
\clearpage

\section*{Ответы и самопроверка}
\vspace{3mm}
\renewcommand{\arraystretch}{1.42}
\begin{center}
\begin{tabularx}{.9\linewidth}{@{}>{\centering\arraybackslash}p{.12\linewidth}X@{}}
\toprule
\textbf{№} & \textbf{Правильный ответ}\\
\midrule
1 & $50$\\
2 & $14$\\
3 & $24$\\
4 & $20$\\
5 & $12$\\
6 & $30$\\
7 & $1:3$\\
8 & $10$\\
9 & $30$\\
10 & $20$\\
\bottomrule
\end{tabularx}
\end{center}
\vspace{6mm}
\textbf{Проверьте метод, если ответ не совпал.}
\begin{itemize}
\item В задачах 3 и 6 важна половина высоты трапеции.
\item В задачах 4 и 5 важно различать отношение отрезков и отношение площади к площади.
\item В задачах 7 и 8 полезно продлить $BC$ до пересечения с $AR$.
\item В задачах 9 и 10 удобно провести среднюю линию трапеции, а затем параллельную ей через $N$.
\end{itemize}
\vfill
\begin{center}
{\footnotesize\color{muted}Лёвин Артём Александрович --- эксклюзивно для Софии Кальней}
\end{center}
\end{document}
